% =============================================================================
%  Chapter 12 — Domain Vertex Transition
% =============================================================================

\chapter{Domain Vertex Transition}
\label{ch:domain-vertex-transition}

Chapter~\ref{ch:domain-ontology} fixed the ontology of the refactor.  This
chapter documents the next implementation step: the domain now treats
\texttt{Vertex} as the canonical geometric storage, while \texttt{Node}
survives as a transitional analysis-layer cache.

The design goal of this step is deliberately modest:

\begin{quote}
move the ownership boundary of \texttt{Domain} toward pure geometry without
breaking the current element, model, and PETSc assembly pipelines.
\end{quote}

That means this is not yet the final state of the refactor.  It is the
bridge state that keeps existing kernels working while removing the strongest
geometric/analysis conflation.


% ─────────────────────────────────────────────────────────────────────────
\section{Motivation}
\label{sec:domain-vertex-transition-motivation}

Before this step, \texttt{Domain} owned a vector of \texttt{Node<dim>}.  That
had three consequences:

\begin{enumerate}[nosep]
  \item the geometric domain depended directly on nodal DOF storage;
  \item simply constructing a domain implied that the domain was already in an
        analysis-ready state;
  \item purely geometric consumers such as mesh export were forced to read
        geometry through analysis nodes.
\end{enumerate}

This was workable, but it obscured the intended layering:

\[
\texttt{Domain} \not\equiv \texttt{Model}.
\]

The refactor therefore changes the ownership model to:

\[
\texttt{Domain owns vertices}, \qquad
\texttt{Model / assembly consume nodes}.
\]


% ─────────────────────────────────────────────────────────────────────────
\section{Implemented Transition Design}
\label{sec:domain-vertex-transition-design}

\subsection{Canonical storage}

\texttt{Domain} now stores
\[
\texttt{std::vector<Vertex<dim>>}
\]
as its primary geometric container.  This is the authoritative representation
of mesh points for:

\begin{itemize}[nosep]
  \item coordinates;
  \item mesh identity;
  \item sorting and lookup by vertex ID;
  \item geometry-only consumers.
\end{itemize}

\subsection{Node cache}

To preserve compatibility with the current element and model infrastructure,
\texttt{Domain} also materialises a
\[
\texttt{std::vector<Node<dim>>}
\]
cache on demand.  This cache is derived from the vertices and is used only when
analysis-level operations require nodal DOFs or PETSc indexing.

This means the runtime interpretation is now:

\begin{itemize}[nosep]
  \item \texttt{vertices\_} is the geometric truth;
  \item \texttt{nodes\_cache\_} is a derived analysis view.
\end{itemize}

\paragraph{Why this is acceptable.}
The cache is not introduced for performance in the hot path.  It is a
compatibility layer during migration.  Since nodes are created once and then
reused during analysis, the cost is negligible relative to element assembly and
material updates.

\subsection{Cache invalidation rule}

Whenever the vertex container changes, the node cache is invalidated.  The next
request for \texttt{nodes()}, \texttt{node(id)}, or element-node binding
rebuilds the cache.

This makes the design explicit:

\begin{quote}
geometry changes invalidate analysis views.
\end{quote}

That is the correct dependency direction.


% ─────────────────────────────────────────────────────────────────────────
\section{Why This Preserves Performance}
\label{sec:domain-vertex-transition-performance}

The user requirement for this refactor was clear: do not sacrifice runtime
performance unnecessarily, and favour compile-time structure suitable for HPC.

This implementation respects that requirement in three ways.

\subsection{No change to element hot loops}

The integration, constitutive update, and local stiffness computations remain
exactly where they already were:

\begin{itemize}[nosep]
  \item \texttt{ElementGeometry};
  \item \texttt{ContinuumElement};
  \item \texttt{BeamElement};
  \item \texttt{TimoshenkoBeamN};
  \item \texttt{ShellElement}.
\end{itemize}

Those kernels still operate on bound \texttt{Node} objects and continue to
benefit from fixed-size matrices, compile-time quadrature, and policy-based
dispatch.

\subsection{Geometry-only consumers stop depending on nodes}

The VTK mesh-export path was updated to read coordinates from
\texttt{Domain::vertices()} rather than from \texttt{Domain::nodes()}.
This is architecturally cleaner and slightly cheaper conceptually: geometry is
now consumed as geometry.

\subsection{PETSc creation is now lazy}

An important side effect of the refactor is the change to \texttt{Mesh}.  The
\texttt{DM} object is no longer created in the \texttt{Mesh} constructor.
Instead, it is created lazily when the DMPlex-specific operations are first
needed.

This has two benefits:

\begin{enumerate}[nosep]
  \item a purely geometric \texttt{Domain} can be created and manipulated
        without immediately entering PETSc runtime semantics;
  \item tests that exercise only the geometric/domain side no longer require
        \texttt{PetscInitialize}.
\end{enumerate}

Architecturally, this is important because it moves PETSc from
``always present'' to ``activated when the backend is actually used''.


% ─────────────────────────────────────────────────────────────────────────
\section{Code-Level Changes}
\label{sec:domain-vertex-transition-code}

The main code changes in this phase are:

\begin{enumerate}[nosep]
  \item \texttt{Vertex<Dim>} introduced earlier becomes the canonical
        storage type inside \texttt{Domain}.
  \item \texttt{Domain} gains explicit vertex accessors:
        \texttt{vertices()}, \texttt{vertex(id)}, \texttt{num\_vertices()}.
  \item \texttt{Domain::nodes()} becomes a lazily materialised analysis view.
  \item \texttt{GmshDomainBuilder} now sorts the domain through
        \texttt{sort\_vertices\_by\_id()} rather than sorting the nodal view.
  \item \texttt{VTKModelExporter} and the legacy \texttt{VTKDataContainer}
        now use vertices for mesh coordinates.
  \item \texttt{Mesh} now creates the PETSc \texttt{DM} lazily.
\end{enumerate}

This is intentionally a minimal-surface-area change.  The aim was to shift the
ownership boundary without forcing an immediate rewrite of
\texttt{ElementGeometry} or \texttt{Model}.


% ─────────────────────────────────────────────────────────────────────────
\section{What Remains Transitional}
\label{sec:domain-vertex-transition-transitional}

This phase does \emph{not} yet complete the domain refactor.

The following are still transitional:

\begin{itemize}[nosep]
  \item \texttt{ElementGeometry} still binds concrete \texttt{Node} pointers,
        not geometry-only vertex handles.
  \item \texttt{Model} still reads and writes nodal DOF data through the
        domain-level node view.
  \item PETSc sieve IDs are still stored in \texttt{Node}, not yet in a
        dedicated analysis/indexing layer.
\end{itemize}

That is acceptable at this stage.  The important structural improvement is that
these analysis concerns are no longer the \emph{canonical ownership model} of
the domain itself.


% ─────────────────────────────────────────────────────────────────────────
\section{Validation}
\label{sec:domain-vertex-transition-validation}

This phase adds a dedicated test executable for the transition.

\subsection{Domain/Vertex refactor test}

The new test checks:

\begin{enumerate}[nosep]
  \item that vertices are the canonical stored geometry;
  \item that the nodal view is materialised from those vertices;
  \item that sorting vertices updates the nodal view coherently;
  \item that analysis-only data written into the node cache does not alter the
        geometric meaning of the corresponding vertex;
  \item that the legacy \texttt{add\_node(Node\&\&)} path still preserves
        geometry correctly.
\end{enumerate}

\subsection{Regression coverage}

The existing serendipity test suite was recompiled and rerun after this change,
which is a useful regression check because it exercises:

\begin{itemize}[nosep]
  \item \texttt{ElementGeometry};
  \item quadrature;
  \item VTK export;
  \item geometry binding through nodes.
\end{itemize}

In addition, the Gmsh pipeline target was rebuilt successfully, confirming that
the builder and domain interface remain consistent at compile time.


% ─────────────────────────────────────────────────────────────────────────
\section{Architectural Consequences}
\label{sec:domain-vertex-transition-consequences}

This transition has one especially important architectural consequence:

\begin{quote}
\texttt{Domain} can now be understood as a geometric object that happens to be
able to provide an analysis view, rather than as an analysis object that also
stores geometry.
\end{quote}

That inversion is subtle, but it matters.  It clarifies the dependency
direction for the next phase:

\begin{enumerate}[nosep]
  \item first, geometry is owned by the domain;
  \item second, analysis nodes are attached or materialised from it;
  \item third, elements and solvers consume those analysis nodes.
\end{enumerate}

This is much closer to the ontology laid out in the previous chapter.


% ─────────────────────────────────────────────────────────────────────────
\section{Next Step}
\label{sec:domain-vertex-transition-next}

The next logical step is now well defined:

\begin{quote}
remove the direct \texttt{Node}-dependence from \texttt{ElementGeometry}.
\end{quote}

That will require introducing a geometry-facing vertex handle or concept for
binding coordinates, while the model layer separately manages nodal DOFs and
PETSc indexing.

Once that is done, \texttt{Domain} will no longer need to expose a nodal cache
at all, except perhaps as a legacy adapter.


% ─────────────────────────────────────────────────────────────────────────
\section{Conclusion}
\label{sec:domain-vertex-transition-conclusion}

This chapter marks the first real implementation step of the domain refactor.
The main result is not a new feature but a new ownership boundary:

\begin{itemize}[nosep]
  \item vertices own domain geometry;
  \item nodes are an analysis view;
  \item mesh export reads geometry from vertices;
  \item PETSc mesh backend creation is deferred until needed.
\end{itemize}

No element hot path was slowed down, no quadrature logic was generalized away,
and existing element kernels continue to operate unchanged.  That is the right
kind of progress for this refactor: architectural movement without numerical
regression.
